% ============================================================
% PASTE THIS INSIDE \section{Experimental Setup}
% Replace the existing single paragraph that begins:
%   "We evaluate variable-role representations in Qwen2.5-1.5B..."
% ============================================================

We evaluate variable-role representations in \textbf{Qwen2.5-1.5B} with frozen
weights. Our data combines realistic functions from \textbf{CodeSearchNet}
\citep{husain2019codesearchnet} with aligned multilingual code from
\textbf{XLCoST} \citep{zhu2022xlcost} and \textbf{MuST-CoST}
\citep{zhu2022mustcost}. We label variable occurrences by role---accumulator,
index, iterator, boolean, and class/structure reference---using lightweight
static analysis plus a small manual audit, and add controlled renaming and
role-swapping counterfactuals. We train linear probes on hidden states and
report per-role and macro-averaged F1 on held-out programs, renamed variants,
and cross-dataset examples.

Each role captures a distinct functional pattern in how a variable is used
within a program. An \textbf{accumulator} is a variable that aggregates values
across iterations of a loop, typically updated via \texttt{+=}, \texttt{-=},
\texttt{.append()}, or \texttt{.extend()} inside a loop body (e.g.,
\texttt{total}, \texttt{count}, \texttt{result}). An \textbf{index} (or key)
variable serves as a loop counter or dictionary key, appearing in \texttt{for}-loop
headers or subscript expressions (e.g., \texttt{i}, \texttt{j}, \texttt{idx}).
An \textbf{iterator} is a variable bound in a \texttt{for}-loop header that
ranges over a sequence; we track all subsequent uses of that variable within
the loop body via scope-aware AST traversal. A \textbf{boolean control}
variable is used as a conditional guard---in \texttt{if} tests, boolean
assignments, loop conditions, return expressions, or indexing contexts---rather
than holding numeric or string data; each occurrence is labeled by its usage
context (e.g., \texttt{conditional\_use}, \texttt{loop\_use},
\texttt{return\_use}). A \textbf{class/structure reference} is a variable
whose type is a declared class or struct; we locate each \texttt{class}/\texttt{struct}
definition via brace matching (in C++ and C\#), then label the type name at its
declaration site and every subsequent occurrence of that name in the snippet.
These five roles span the range from structurally-driven (accumulator, iterator)
to lexically-driven (index) representations, and from single-language
(boolean, class reference) to cross-language (accumulator, index) detection
patterns---motivating the different probing methodologies described in
Section~\ref{sec:evidence}.

\paragraph{Renaming scope and implementations.}
Renaming counterfactuals are evaluated on \textbf{Python only}; the
seven-language corpus supports baseline probing and cross-language
transfer, not renaming (the pipeline refuses renaming requests for other
languages rather than emitting corrupted source, since whole-word
substitution rewrites keywords, preprocessor directives, and stdlib names
outside Python). Two renaming implementations appear in this work and we
report which produced each result: the pipeline's whole-word substitution
(index, accumulator, iterator, and class/structure results), which
rewrites every textual occurrence of a target name including inside
string literals and comments, and the boolean workstream's scope-aware
span editor (boolean results), which renames only identifier occurrences,
preserves strings and comments, and re-verifies each rewritten program by
parsing and occurrence-identity checks. Both preserve program parseability
on Python; converging on a single parser-based renamer is planned
follow-up work.
